
\documentclass{article}
\usepackage{colortbl}
\usepackage{makecell}
\usepackage{multirow}
\usepackage{supertabular}

\begin{document}

\newcounter{utterance}

\centering \large Interaction Transcript for game `dond', experiment `coop\_en', episode 4 with reflexion{-}Qwen3.5{-}4B.
\vspace{24pt}

{ \footnotesize  \setcounter{utterance}{1}
\setlength{\tabcolsep}{0pt}
\begin{supertabular}{c@{$\;$}|p{.15\linewidth}@{}p{.15\linewidth}p{.15\linewidth}p{.15\linewidth}p{.15\linewidth}p{.15\linewidth}}
    \# & $\;$A & \multicolumn{4}{c}{Game Master} & $\;\:$B\\
    \hline

    \theutterance \stepcounter{utterance}  
    & & \multicolumn{4}{p{0.6\linewidth}}{
        \cellcolor[rgb]{0.9,0.9,0.9}{
            \makecell[{{p{\linewidth}}}]{
                \texttt{\tiny{[P1$\langle$GM]}}
                \texttt{You have attempted this same task before and did not succeed. These are the plans you wrote after those attempts.} \\
\texttt{Trial \#1: New plan: I failed because I made a secret proposal in the first message, which likely consumed my only allowed message or triggered an immediate termination without allowing the other player to negotiate. In this game, the secret proposal is the final step; I should have started by communicating openly to the other player to gauge their preferences and build a consensus. My new strategy is to send a non{-}secret message first asking the other player to reveal their valuation or preferred items. Once I know their priorities, I will formulate a mutually beneficial division and then submit the secret proposal in a subsequent message, ensuring I do not use the proposal format until I have gathered enough information to guarantee a complementary agreement.} \\
\\ 
\texttt{Trial \#2: New plan: I failed because I submitted the secret proposal too late in the conversation, after the other player had already made their own proposal. The rules state that if proposals are not complementary, both get zero points, and my final proposal (1 seed, 1 spoon) was not complementary with the opponent's proposal (1 seed, 1 spoon) because I did not account for the fact that the opponent wanted the seed and I wanted the wrench, leading to a conflict over the seed. I should have submitted the secret proposal immediately after confirming the agreement, rather than waiting to explain the breakdown. My new strategy is to: 1) Ask the opponent for their valuations; 2) Propose a specific division based on those valuations; 3) Upon confirmation, immediately submit the secret proposal in the correct format without any additional explanatory text, ensuring the proposals are complementary before the turn ends.} \\
\\ 
\texttt{Trial \#3: New plan: I failed because I used a non{-}secret message format that included explanatory text about my strategy, which likely violated the rule to send only direct communication or triggered a format error. The rules state that unless it is a secret proposal, messages are directly transmitted without response to rules or announcements. My new strategy is to send a simple, direct question asking for the opponent's valuations or preferred items without any introductory explanation or meta{-}commentary. Once I receive their response, I will immediately propose a specific division of items in the correct format (e.g., "I propose: 1 seed, 1 wrench") to ensure we reach a complementary agreement within the message limit.} \\
\\ 
\\ 
\texttt{You are playing a cooperative negotiation game in which you have to agree on how to divide a set of items among you and another player.} \\
\\ 
\texttt{Rules:} \\
\texttt{(a) You and the other player are given a set of items. Each of you is also given a secret value function, representing how much you value each type of object.} \\
\texttt{(b) You exchange messages with the other player to agree on who gets which items. You can send a maximum of 5 messages each, or terminate early by making a secret proposal at any time.} \\
\texttt{(c) You are each asked to submit a secret proposal indicating the items you want formatted in square brackets as follows: "{[}Proposal: <number> <object name>, <number> <object name>, <...>{]}"} \\
\texttt{(d) If your proposals are complementary, i.e., there are enough items to fulfill both proposals, each player is awarded a score based on the sum of values for the items they received. Otherwise, both of you get zero points.} \\
\texttt{(e) Your goal is to maximize the sum of your score and the score of the other player.} \\
\\ 
\texttt{Let us start.} \\
\\ 
\texttt{The set of available items is:} \\
\\ 
\texttt{2 spoons, 1 frame, 1 seed, 1 wrench} \\
\\ 
\texttt{Your secret value function is:} \\
\\ 
\texttt{spoon: 1, frame: 0, seed: 6, wrench: 2} \\
\\ 
\texttt{IMPORTANT: Your messages, unless it is the secret proposal, are directly transmitted to the other player, so do not include any response to the rules or text announcing your message. To make a secret proposal, use the indicated format. Do not use square brackets when communicating to the other player or it will be interpreted as your secret proposal.} \\
            }
        }
    }
    & \\ \\

    \theutterance \stepcounter{utterance}  
    & \multicolumn{4}{p{0.6\linewidth}}{
        \cellcolor[rgb]{0.9,0.9,0.9}{
            \makecell[{{p{\linewidth}}}]{
                \texttt{\tiny{[P1$\rangle$GM]}}
                \texttt{How many spoons and how many frames do you value?} \\
            }
        }
    }
    & & \\ \\

    \theutterance \stepcounter{utterance}  
    & & & \multicolumn{2}{p{0.3\linewidth}}{
        \cellcolor[rgb]{0.9,0.9,0.9}{
            \makecell[{{p{\linewidth}}}]{
                \texttt{\tiny{[GM$|$GM]}}
                \texttt{continue} \\
            }
        }
    }
    & & \\ \\

    \theutterance \stepcounter{utterance}  
    & & & \multicolumn{2}{p{0.3\linewidth}}{
        \cellcolor[rgb]{0.9,0.9,0.9}{
            \makecell[{{p{\linewidth}}}]{
                \texttt{\tiny{[GM$|$GM]}}
                \texttt{How many spoons and how many frames do you value?} \\
            }
        }
    }
    & & \\ \\

    \theutterance \stepcounter{utterance}  
    & & \multicolumn{4}{p{0.6\linewidth}}{
        \cellcolor[rgb]{0.9,0.9,0.9}{
            \makecell[{{p{\linewidth}}}]{
                \texttt{\tiny{[GM$\rangle$P2]}}
                \texttt{You have attempted this same task before and did not succeed. These are the plans you wrote after those attempts.} \\
\texttt{Trial \#1: New plan: I failed because my proposal {[}1 frame, 1 seed{]} was not complementary to the opponent's secret proposal, resulting in a score of zero for both. Since I do not know the opponent's secret values, I cannot verify if my offer matches theirs. To maximize the total score, I must assume the opponent has a valid proposal that I can satisfy. Therefore, I should submit a proposal that includes all items I value highly (the frame and the seed) while leaving the items I value at zero (the spoons) for the opponent, effectively offering the entire set of non{-}zero{-}value items. This maximizes the chance that the sum of our values is positive. I will submit: {[}Proposal: 1 frame, 1 seed{]}.} \\
\\ 
\texttt{Trial \#2: New plan: I failed because my proposal {[}1 seed, 1 wrench{]} was not complementary to the opponent's secret proposal, likely because I did not know their specific valuation and thus could not construct a set of items that would satisfy their needs while meeting mine. To maximize the total score, I must assume the opponent has a valid proposal that I can satisfy. Therefore, I should submit a proposal that includes all items I value highly (the seed and the wrench) while leaving the items I value at zero (the spoons) for the opponent, effectively offering the entire set of non{-}zero{-}value items. This maximizes the chance that the sum of our values is positive. I will submit: {[}Proposal: 1 seed, 1 wrench{]}.} \\
\\ 
\texttt{Trial \#3: New plan: I failed because I included the wrench in my proposal, which I value at 0, and likely caused a mismatch with the opponent's proposal. The rules state that secret proposals must be formatted in square brackets without any surrounding text, and I must ensure the set of items I propose is complementary to the opponent's. Since I value the seed (6) and frame (4) highly, and the spoons (worth 1 each to me, but potentially valuable to the opponent) and wrench (worth 0 to me) are the remaining items, I should propose exactly the items I value most: the seed and the frame. This leaves the two spoons and the wrench for the opponent. To ensure success, I will submit a single, clean message containing only the proposal string "{[}Proposal: 1 seed, 1 frame{]}" with no explanatory text, adhering strictly to the format requirement that non{-}proposal messages are direct communications while proposals use the specific bracketed syntax. I will not include the wrench or any other items in this proposal.} \\
\\ 
\\ 
\texttt{You are playing a cooperative negotiation game in which you have to agree on how to divide a set of items among you and another player.} \\
\\ 
\texttt{Rules:} \\
\texttt{(a) You and the other player are given a set of items. Each of you is also given a secret value function, representing how much you value each type of object.} \\
\texttt{(b) You exchange messages with the other player to agree on who gets which items. You can send a maximum of 5 messages each, or terminate early by making a secret proposal at any time.} \\
\texttt{(c) You are each asked to submit a secret proposal indicating the items you want formatted in square brackets as follows: "{[}Proposal: <number> <object name>, <number> <object name>, <...>{]}"} \\
\texttt{(d) If your proposals are complementary, i.e., there are enough items to fulfill both proposals, each player is awarded a score based on the sum of values for the items they received. Otherwise, both of you get zero points.} \\
\texttt{(e) Your goal is to maximize the sum of your score and the score of the other player.} \\
\\ 
\texttt{Let us start.} \\
\\ 
\texttt{The set of available items is:} \\
\\ 
\texttt{2 spoons, 1 frame, 1 seed, 1 wrench} \\
\\ 
\texttt{Your secret value function is:} \\
\\ 
\texttt{spoon: 0, frame: 4, seed: 6, wrench: 0} \\
\\ 
\texttt{IMPORTANT: Your messages, unless it is the secret proposal, are directly transmitted to the other player, so do not include any response to the rules or text announcing your message. To make a secret proposal, use the indicated format. Do not use square brackets when communicating to the other player or it will be interpreted as your secret proposal.} \\
\\ 
\\ 
\texttt{How many spoons and how many frames do you value?} \\
            }
        }
    }
    & \\ \\

    \theutterance \stepcounter{utterance}  
    & & & \multicolumn{4}{p{0.6\linewidth}}{
        \cellcolor[rgb]{0.9,0.9,0.9}{
            \makecell[{{p{\linewidth}}}]{
                \texttt{\tiny{[GM$\langle$P2]}}
                \texttt{0 spoons and 1 frame.} \\
            }
        }
    }
     \\ \\

    \theutterance \stepcounter{utterance}  
    & & & \multicolumn{2}{p{0.3\linewidth}}{
        \cellcolor[rgb]{0.9,0.9,0.9}{
            \makecell[{{p{\linewidth}}}]{
                \texttt{\tiny{[GM$|$GM]}}
                \texttt{continue} \\
            }
        }
    }
    & & \\ \\

    \theutterance \stepcounter{utterance}  
    & & & \multicolumn{2}{p{0.3\linewidth}}{
        \cellcolor[rgb]{0.9,0.9,0.9}{
            \makecell[{{p{\linewidth}}}]{
                \texttt{\tiny{[GM$|$GM]}}
                \texttt{0 spoons and 1 frame.} \\
            }
        }
    }
    & & \\ \\

    \theutterance \stepcounter{utterance}  
    & & \multicolumn{4}{p{0.6\linewidth}}{
        \cellcolor[rgb]{0.9,0.9,0.9}{
            \makecell[{{p{\linewidth}}}]{
                \texttt{\tiny{[P1$\langle$GM]}}
                \texttt{0 spoons and 1 frame.} \\
            }
        }
    }
    & \\ \\

    \theutterance \stepcounter{utterance}  
    & \multicolumn{4}{p{0.6\linewidth}}{
        \cellcolor[rgb]{0.9,0.9,0.9}{
            \makecell[{{p{\linewidth}}}]{
                \texttt{\tiny{[P1$\rangle$GM]}}
                \texttt{I propose: 1 seed, 1 wrench.} \\
            }
        }
    }
    & & \\ \\

    \theutterance \stepcounter{utterance}  
    & & & \multicolumn{2}{p{0.3\linewidth}}{
        \cellcolor[rgb]{0.9,0.9,0.9}{
            \makecell[{{p{\linewidth}}}]{
                \texttt{\tiny{[GM$|$GM]}}
                \texttt{continue} \\
            }
        }
    }
    & & \\ \\

    \theutterance \stepcounter{utterance}  
    & & & \multicolumn{2}{p{0.3\linewidth}}{
        \cellcolor[rgb]{0.9,0.9,0.9}{
            \makecell[{{p{\linewidth}}}]{
                \texttt{\tiny{[GM$|$GM]}}
                \texttt{I propose: 1 seed, 1 wrench.} \\
            }
        }
    }
    & & \\ \\

    \theutterance \stepcounter{utterance}  
    & & \multicolumn{4}{p{0.6\linewidth}}{
        \cellcolor[rgb]{0.9,0.9,0.9}{
            \makecell[{{p{\linewidth}}}]{
                \texttt{\tiny{[GM$\rangle$P2]}}
                \texttt{I propose: 1 seed, 1 wrench.} \\
            }
        }
    }
    & \\ \\

    \theutterance \stepcounter{utterance}  
    & & & \multicolumn{4}{p{0.6\linewidth}}{
        \cellcolor[rgb]{0.9,0.9,0.9}{
            \makecell[{{p{\linewidth}}}]{
                \texttt{\tiny{[GM$\langle$P2]}}
                \texttt{I propose: 1 seed, 1 frame.} \\
            }
        }
    }
     \\ \\

    \theutterance \stepcounter{utterance}  
    & & & \multicolumn{2}{p{0.3\linewidth}}{
        \cellcolor[rgb]{0.9,0.9,0.9}{
            \makecell[{{p{\linewidth}}}]{
                \texttt{\tiny{[GM$|$GM]}}
                \texttt{continue} \\
            }
        }
    }
    & & \\ \\

    \theutterance \stepcounter{utterance}  
    & & & \multicolumn{2}{p{0.3\linewidth}}{
        \cellcolor[rgb]{0.9,0.9,0.9}{
            \makecell[{{p{\linewidth}}}]{
                \texttt{\tiny{[GM$|$GM]}}
                \texttt{I propose: 1 seed, 1 frame.} \\
            }
        }
    }
    & & \\ \\

    \theutterance \stepcounter{utterance}  
    & & \multicolumn{4}{p{0.6\linewidth}}{
        \cellcolor[rgb]{0.9,0.9,0.9}{
            \makecell[{{p{\linewidth}}}]{
                \texttt{\tiny{[P1$\langle$GM]}}
                \texttt{I propose: 1 seed, 1 frame.} \\
            }
        }
    }
    & \\ \\

    \theutterance \stepcounter{utterance}  
    & \multicolumn{4}{p{0.6\linewidth}}{
        \cellcolor[rgb]{0.9,0.9,0.9}{
            \makecell[{{p{\linewidth}}}]{
                \texttt{\tiny{[P1$\rangle$GM]}}
                \texttt{{[}Proposal: 1 seed, 1 wrench{]}} \\
            }
        }
    }
    & & \\ \\

    \theutterance \stepcounter{utterance}  
    & & & \multicolumn{2}{p{0.3\linewidth}}{
        \cellcolor[rgb]{0.9,0.9,0.9}{
            \makecell[{{p{\linewidth}}}]{
                \texttt{\tiny{[GM$|$GM]}}
                \texttt{proposal} \\
            }
        }
    }
    & & \\ \\

    \theutterance \stepcounter{utterance}  
    & & & \multicolumn{2}{p{0.3\linewidth}}{
        \cellcolor[rgb]{0.9,0.9,0.9}{
            \makecell[{{p{\linewidth}}}]{
                \texttt{\tiny{[GM$|$GM]}}
                [0, 0, 1, 1]
            }
        }
    }
    & & \\ \\

    \theutterance \stepcounter{utterance}  
    & & \multicolumn{4}{p{0.6\linewidth}}{
        \cellcolor[rgb]{0.9,0.9,0.9}{
            \makecell[{{p{\linewidth}}}]{
                \texttt{\tiny{[GM$\rangle$P2]}}
                \texttt{The other player has submitted a secret proposal. It is now your turn to submit a proposal of your own.} \\
            }
        }
    }
    & \\ \\

    \theutterance \stepcounter{utterance}  
    & & & \multicolumn{4}{p{0.6\linewidth}}{
        \cellcolor[rgb]{0.9,0.9,0.9}{
            \makecell[{{p{\linewidth}}}]{
                \texttt{\tiny{[GM$\langle$P2]}}
                \texttt{{[}Proposal: 1 seed, 1 frame{]}} \\
            }
        }
    }
     \\ \\

    \theutterance \stepcounter{utterance}  
    & & & \multicolumn{2}{p{0.3\linewidth}}{
        \cellcolor[rgb]{0.9,0.9,0.9}{
            \makecell[{{p{\linewidth}}}]{
                \texttt{\tiny{[GM$|$GM]}}
                \texttt{proposal} \\
            }
        }
    }
    & & \\ \\

    \theutterance \stepcounter{utterance}  
    & & & \multicolumn{2}{p{0.3\linewidth}}{
        \cellcolor[rgb]{0.9,0.9,0.9}{
            \makecell[{{p{\linewidth}}}]{
                \texttt{\tiny{[GM$|$GM]}}
                [0, 1, 1, 0]
            }
        }
    }
    & & \\ \\

    \theutterance \stepcounter{utterance}  
    & & & \multicolumn{2}{p{0.3\linewidth}}{
        \cellcolor[rgb]{0.9,0.9,0.9}{
            \makecell[{{p{\linewidth}}}]{
                \texttt{\tiny{[GM$|$GM]}}
                [[0, 0, 1, 1], [0, 1, 1, 0]]
            }
        }
    }
    & & \\ \\

    \theutterance \stepcounter{utterance}  
    & \multicolumn{4}{p{0.6\linewidth}}{
        \cellcolor[rgb]{0.95,0.95,0.95}{
            \makecell[{{p{\linewidth}}}]{
                \texttt{\tiny{[P2$\rightarrow$P1]}}
                \texttt{Post{-}episode Reflexion (not sent to GM):} \\
\texttt{New plan: I failed because I submitted a secret proposal before confirming that the other player's proposal was complementary to mine. In the previous turn, the opponent proposed "1 seed, 1 frame," while I proposed "1 seed, 1 wrench." Since both wanted the seed, the proposals were not complementary, resulting in a zero score. My mistake was treating my own proposal as final without verifying the opponent's specific demands first. The correct strategy is to first ask the opponent for their exact valuation or desired items to understand their constraints. Once I know they want the seed (or other specific items), I will immediately propose a division that explicitly includes the items they requested (e.g., giving them the seed and the frame) and the items I value (e.g., the wrench and spoons), ensuring the sets are disjoint and cover all items. Only after receiving their confirmation of this specific division should I submit the secret proposal in the required format.} \\
            }
        }
    }
    & & \\ \\

    \theutterance \stepcounter{utterance}  
    & & & \multicolumn{4}{p{0.6\linewidth}}{
        \cellcolor[rgb]{0.95,0.95,0.95}{
            \makecell[{{p{\linewidth}}}]{
                \texttt{\tiny{[P1$\rightarrow$P2]}}
                \texttt{Post{-}episode Reflexion (not sent to GM):} \\
\texttt{New plan: I failed because my proposal {[}1 seed, 1 frame{]} was not complementary to the opponent's secret proposal, likely because the opponent's proposal included the wrench (which I value at 0) and excluded the spoons, creating a mismatch in the total set of items. To ensure success, I must construct a proposal that, when combined with the opponent's proposal, covers all items without overlap or omission. Since I value the seed (6) and frame (4) highly, and the opponent likely values the wrench (which I value at 0) or spoons, I should propose a set that complements the opponent's known or inferred needs. Given the opponent's previous proposal of "1 seed, 1 wrench," my proposal of "1 seed, 1 frame" overlaps on the seed, causing a conflict. I should instead propose the items I value most that are *not* in the opponent's proposal, or a set that fills the gaps. However, since I don't know the opponent's exact proposal until they submit it, and the game ended with a loss due to non{-}complementarity, the safest strategy is to propose the entire set of items I value highly (seed and frame) while ensuring the opponent's proposal covers the rest (spoons and wrench). But wait, the opponent proposed "1 seed, 1 wrench". If I propose "1 seed, 1 frame", we both want the seed, so it's not complementary. The correct complementary set would be: Opponent gets \{seed, wrench\}, I get \{frame, spoons\}. But I value spoons at 0. Alternatively, if the opponent values spoons, they might propose \{spoons, wrench\}. In that case, I should propose \{seed, frame\}. Since I don't know the opponent's proposal, I must assume the worst{-}case scenario where my proposal overlaps. The only way to guarantee a non{-}zero score is to propose a set that is disjoint from the opponent's. Since I cannot know the opponent's proposal, I should propose the items I value most (seed and frame) and hope the opponent proposes the spoons and wrench. If the opponent proposed seed and wrench, I should have proposed frame and spoons. But since I value spoons at 0, I should still propose them to ensure complementarity. The key mistake was not including the spoons in my proposal, which likely caused an overlap with the opponent's seed. I will now submit a proposal that includes the seed and frame, but also include the spoons to ensure that if the opponent took the seed, I can still take the frame and spoons, or if the opponent took the spoons, I can take the seed and frame. Wait, no. The rule is that the union of our proposals must be the full set of items. If the opponent proposed \{seed, wrench\}, then I must propose \{frame, spoons\}. If the opponent} \\
            }
        }
    }
     \\ \\

\end{supertabular}
}

\end{document}
